\documentclass[11pt,a4paper]{article}

\usepackage[utf8]{inputenc}
\usepackage[margin=1.15in]{geometry}
\usepackage{amsmath,amssymb,amsthm,mathtools}
\usepackage{booktabs,longtable,array}
\usepackage{enumitem}
\usepackage{xcolor}
\usepackage{microtype}
\usepackage[colorlinks=true,linkcolor=blue!50!black,citecolor=green!40!black,urlcolor=blue!50!black]{hyperref}

\setlist{itemsep=2pt,topsep=4pt,parsep=0pt}
\numberwithin{equation}{section}
\allowdisplaybreaks

% ---------- theorem environments ----------
\theoremstyle{plain}
\newtheorem{theorem}{Theorem}[section]
\newtheorem{proposition}[theorem]{Proposition}
\newtheorem{lemma}[theorem]{Lemma}
\newtheorem{corollary}[theorem]{Corollary}
\theoremstyle{definition}
\newtheorem{definition}[theorem]{Definition}
\newtheorem{example}[theorem]{Example}
\theoremstyle{remark}
\newtheorem{remark}[theorem]{Remark}

% ---------- macros ----------
\newcommand{\R}{\mathbb{R}}
\newcommand{\Z}{\mathbb{Z}}
\newcommand{\frakg}{\mathfrak{g}}
\newcommand{\Ad}{\operatorname{Ad}}
\newcommand{\ad}{\operatorname{ad}}
\newcommand{\rank}{\operatorname{rank}}
\newcommand{\pb}[2]{\{#1,#2\}}
\newcommand{\pair}[2]{\langle #1,#2\rangle}
\newcommand{\weq}{\mathrel{\approx}}
\newcommand{\dd}{\mathrm{d}}
\newcommand{\inv}{^{-1}}

\title{\textbf{Two Notions of Reduction in Classical Mechanics:\\[2pt] Constraints versus Symmetries}}
\author{}
\date{}

\begin{document}
\maketitle

\begin{abstract}
The word ``reduction'' denotes two different procedures in classical mechanics.
\emph{Reduction by constraints} confines the motion to a submanifold of configuration or phase space by means of reaction forces that are not part of the Lagrangian;
\emph{reduction by symmetry} uses the conserved momenta of a group of symmetries of the dynamics to eliminate variables.
Both lower the number of degrees of freedom and both can be formulated in Lagrangian, Hamiltonian and Poisson language, yet they differ in logical status, in their effect on the equations of motion, and in the geometry of the reduced space.
We develop the two theories side by side, with proofs of the principal results (the Lagrange--d'Alembert principle for holonomic constraints, Dirac's second-class reduction, Routh reduction, and the Marsden--Weinstein--Meyer theorem), and then compare them point by point.
Our main findings are:
(i) in phase space both are instances of \emph{presymplectic reduction}---restrict the symplectic form to a submanifold $C$ and divide by the kernel of the restricted form---where the kernel is trivial for holonomic (second-class) constraints and consists of the orbits of the isotropy group for a level set of a momentum map;
(ii) the level set $J^{-1}(\mu)$ is cut out by constraints that are first class in the isotropy directions and second class in the remaining directions, so symmetry reduction is a special case of Dirac's theory, with $\dim\frakg_\mu$ first-class and $\dim\mathcal O_\mu$ second-class constraints;
(iii) the decisive dynamical difference is that symmetry reduction projects \emph{genuine} solutions of the unreduced equations (the constraint force vanishes identically), whereas constraint reduction replaces the dynamics by a different one containing reaction forces;
(iv) the two procedures have dual pitfalls: naive substitution of a conserved momentum into the Lagrangian is wrong, and so is naive restriction of the Hamiltonian to a constraint surface;
(v) nonholonomic constraints fall outside the common symplectic framework.
The plane and spherical pendulum, the Kepler problem and the rigid body are worked out as illustrations.
\end{abstract}

\tableofcontents

%======================================================================
\section{Introduction}\label{sec:intro}
%======================================================================

The word \emph{reduction} is used in classical mechanics for two procedures that are routinely confused.
They have the same effect on the bookkeeping---each lowers the number of degrees of freedom---and in the simplest examples both are carried out by the same elementary manoeuvre of ``solving an equation and eliminating a variable''.

\begin{description}
\item[Reduction by constraints.] As part of the definition of the model, the motion is required to take place on a submanifold $C$ of configuration space or phase space: a bead on a wire, a pendulum bob at a fixed distance from its pivot, the particles of a rigid body at fixed mutual distances, a skate that cannot slip sideways. Reaction forces that are not part of the Lagrangian enforce the restriction.
\item[Reduction by symmetry.] The Lagrangian or Hamiltonian is invariant under a group $G$ of transformations. By Noether's theorem there are conserved momenta; fixing their values and discarding the corresponding group variables produces a smaller system: the radial problem for a particle in a central field, Euler's equations for the free rigid body.
\end{description}

Both procedures are classical and each has a canonical literature.
For constraints: Lagrange and d'Alembert \cite{dalembert,lagrange}, Bergmann and Dirac \cite{bergmann1949,dirac1950,dirac1964}, the geometric treatment of Gotay, Nester and Hinds \cite{gnh1978}, and, for velocity-dependent constraints, the nonholonomic literature \cite{akn,bloch2003}.
For symmetry: Routh \cite{routh1877}, Poincar\'e \cite{poincare1901}, Smale \cite{smale1970}, Meyer \cite{meyer1973}, Marsden and Weinstein \cite{mw1974}, and the monographs \cite{am1978,marsden1992,mr1999,op2004}.
The two literatures use different vocabulary and are usually taught in different chapters, and one finds statements such as ``a symmetry is a kind of constraint'' or ``a constraint is a kind of symmetry''.
The purpose of this paper is to say precisely in which senses such statements are true, false, or true only after qualification.

We ask three questions.
\begin{enumerate}[label=(Q\arabic*)]
\item What is reduced, and by what operation on the underlying geometric structure (the Lagrangian, the symplectic form, the Poisson bracket)?
\item Are the trajectories of the reduced system trajectories of the original system? That is, is reduction a \emph{consequence} of the dynamics or a \emph{modification} of it?
\item Is one procedure a special case of the other, and if not, where do they meet?
\end{enumerate}

\paragraph{Summary of results.}
\begin{enumerate}[label=(R\arabic*)]
\item \emph{Common form.} For a symplectic manifold $(M,\omega)$ and a submanifold $C\subset M$ for which $\iota_C^*\omega$ has constant rank, the leaf space of its kernel is symplectic (presymplectic reduction, Theorem~\ref{thm:presymp}). Holonomic constraints give $C$ symplectic (kernel $0$, ``pure restriction''); a level set $J^{-1}(\mu)$ of a momentum map gives a kernel equal to the orbits of the isotropy group $G_\mu$ (``restriction \emph{and} quotient'').
\item \emph{Dynamics.} The equation of motion on $C$ is the same in both cases, $\iota_X\iota_C^*\omega=\dd(H|_C)$. Its solution is $X=X_H+F_C$ where $F_C$ is a ``constraint force''. For symmetry reduction $F_C=0$ (the dynamics itself preserves $C$); for holonomic constraints $F_C\neq0$ in general (Theorem~\ref{thm:samesolve}).
\item \emph{Dirac classification.} The constraints $J-\mu=0$ split into $\dim G_\mu$ first-class and $\dim\mathcal O_\mu$ second-class ones; for abelian $G$ all are first class (Proposition~\ref{prop:symclass}).
\item \emph{Duality of pitfalls.} In the Lagrangian picture constraint reduction is a trivial restriction while symmetry reduction requires the Routhian; in the Hamiltonian picture the roles are reversed (Section~\ref{sec:compare}).
\item \emph{Limits.} Nonholonomic constraints do not fit the symplectic framework; symmetry reduction can fail to produce a manifold (singular reduction); constraint reduction fails when the constraints are dependent.
\end{enumerate}

\paragraph{Outline.}
Section~\ref{sec:prelim} fixes notation. Sections~\ref{sec:constraints} and~\ref{sec:symmetry} develop constraint reduction and symmetry reduction independently. Section~\ref{sec:unified} places both in the presymplectic framework. Section~\ref{sec:examples} contains four worked examples. Section~\ref{sec:compare} is the systematic comparison, including a summary table. Section~\ref{sec:interplay} discusses gauge systems, time reparametrization, and nonholonomic systems with symmetry.

%======================================================================
\section{Preliminaries and conventions}\label{sec:prelim}
%======================================================================

\subsection{Lagrangian and Hamiltonian mechanics}

Let $Q$ be a smooth $n$-manifold with local coordinates $q=(q^1,\dots,q^n)$ and tangent bundle $TQ$. A (time-independent) Lagrangian is a smooth function $L:TQ\to\R$ and its extremals satisfy the Euler--Lagrange equations
\begin{equation}\label{eq:EL}
\frac{\dd}{\dd t}\frac{\partial L}{\partial \dot q^{i}}-\frac{\partial L}{\partial q^{i}}=0 .
\end{equation}
Most of the paper concerns \emph{mechanical} Lagrangians
\begin{equation}\label{eq:mechL}
L(q,\dot q)=\tfrac12 g_{ij}(q)\dot q^i\dot q^j-V(q),
\end{equation}
with $g$ a Riemannian metric (kinetic energy) and $V$ a potential. The Legendre transformation $p_i=\partial L/\partial\dot q^i=g_{ij}\dot q^j$ is a diffeomorphism $TQ\to T^*Q$ and the energy becomes
\begin{equation}\label{eq:mechH}
H(q,p)=\tfrac12 g^{ij}(q)p_ip_j+V(q).
\end{equation}

The cotangent bundle carries the canonical one-form $\theta=p_i\,\dd q^i$ and the symplectic form
\begin{equation}\label{eq:omega}
\omega=-\dd\theta=\dd q^i\wedge\dd p_i .
\end{equation}
For a symplectic manifold $(M,\omega)$ and $F\in C^\infty(M)$ the Hamiltonian vector field $X_F$ is defined by
\begin{equation}\label{eq:XH}
\iota_{X_F}\omega=\dd F,
\end{equation}
so that for $M=T^*Q$ the integral curves of $X_H$ are the solutions of $\dot q=\partial H/\partial p$, $\dot p=-\partial H/\partial q$. The Poisson bracket is
\begin{equation}\label{eq:PB}
\pb FG=\omega(X_F,X_G)=\frac{\partial F}{\partial q^i}\frac{\partial G}{\partial p_i}-\frac{\partial F}{\partial p_i}\frac{\partial G}{\partial q^i},
\end{equation}
so that $\dot F=\pb FH$ along the flow of $X_H$. It satisfies the Jacobi identity, equivalent to $X_{\pb FG}=-[X_F,X_G]$ (see \cite{am1978,mr1999}). The pairing of $\frakg^*$ with $\frakg$ is denoted $\pair\cdot\cdot$.

\subsection{Group actions, momentum maps and Noether's theorem}\label{sec:momap}

Let a Lie group $G$ with Lie algebra $\frakg$ act on a symplectic manifold $(M,\omega)$ by symplectic maps $\Phi_g$. For $\xi\in\frakg$ let $\xi_M(z)=\frac{\dd}{\dd t}\big|_{0}\Phi_{\exp t\xi}(z)$ be the infinitesimal generator. The action is \emph{Hamiltonian} if there is a map $J:M\to\frakg^*$, the \emph{momentum map}, with
\begin{equation}\label{eq:momap}
\iota_{\xi_M}\omega=\dd J_\xi,\qquad J_\xi(z):=\pair{J(z)}{\xi},
\end{equation}
i.e.\ $\xi_M=X_{J_\xi}$. The momentum map is \emph{equivariant} if $J(\Phi_g z)=\Ad^*_g J(z)$, where $\Ad^*_g\mu$ is the coadjoint action, $\pair{\Ad^*_g\mu}{\eta}=\pair\mu{\Ad_{g\inv}\eta}$. Differentiating, and writing $\pair{\ad^*_\xi\mu}{\eta}:=\pair\mu{[\xi,\eta]}$,
\begin{equation}\label{eq:TJ}
T_zJ\,(\xi_M(z))=-\ad^*_\xi J(z),\qquad
\pb{J_\xi}{J_\eta}=J_{[\xi,\eta]} .
\end{equation}
(The second identity uses that $\xi\mapsto\xi_M$ is a Lie algebra anti-homomorphism for left actions and $X_{\pb FG}=-[X_F,X_G]$.) In a basis $e_a$ of $\frakg$ with $[e_a,e_b]=c^c_{ab}e_c$ we write $J_a=J_{e_a}$, so $\pb{J_a}{J_b}=c^c_{ab}J_c$.

For $\mu\in\frakg^*$ let $G_\mu=\{g:\Ad^*_g\mu=\mu\}$ be the isotropy (stabilizer) group, with Lie algebra
\begin{equation}\label{eq:gmu}
\frakg_\mu=\{\xi\in\frakg:\ \ad^*_\xi\mu=0\}=\{\xi:\ \pair\mu{[\xi,\eta]}=0\ \ \forall\eta\in\frakg\},
\end{equation}
and let $\mathcal O_\mu=G\cdot\mu\cong G/G_\mu$ be the coadjoint orbit, so $\dim\mathcal O_\mu=\dim G-\dim G_\mu$. If $G$ is abelian, $G_\mu=G$ and $\mathcal O_\mu=\{\mu\}$.

\medskip\noindent
\emph{Cotangent lifts.} If $G$ acts on $Q$, the lifted action on $T^*Q$ is Hamiltonian with the equivariant momentum map
\begin{equation}\label{eq:Jcot}
J_\xi(q,p)=\pair{p}{\xi_Q(q)}=p_i\,\xi_Q^i(q).
\end{equation}

\begin{proposition}[Noether]\label{prop:noether}
Let $H\in C^\infty(M)$ be $G$-invariant. Then $J$ is constant along the flow of $X_H$. In Lagrangian terms, if $L:TQ\to\R$ is invariant under the tangent lift of the $G$-action on $Q$, then $\pair{\partial L/\partial\dot q}{\xi_Q}$ is constant along solutions of \eqref{eq:EL} for each $\xi\in\frakg$.
\end{proposition}
\begin{proof}
$\dd J_\xi(X_H)=\omega(\xi_M,X_H)=-\omega(X_H,\xi_M)=-\dd H(\xi_M)=0$ by invariance. The Lagrangian statement follows by transporting the flow with the Legendre map and \eqref{eq:Jcot}.
\end{proof}

\subsection{Submanifolds of a symplectic manifold; first and second class}\label{sec:subman}

For a subspace $W$ of a symplectic vector space $(V,\Omega)$ put $W^\Omega=\{v:\Omega(v,w)=0\ \forall w\in W\}$; then $\dim W^\Omega=\dim V-\dim W$ and the kernel of $\Omega|_W$ is $W\cap W^\Omega$. $W$ is \emph{symplectic} if $W\cap W^\Omega=0$, \emph{coisotropic} if $W^\Omega\subset W$ and \emph{Lagrangian} if $W=W^\Omega$. A submanifold $C\subset M$ is symplectic, coisotropic, \dots\ if every $T_zC$ is.

Let $C=\{z:\chi_A(z)=0,\ A=1,\dots,m\}$ with $\dd\chi_A$ independent on $C$ and put $\Delta_{AB}:=\pb{\chi_A}{\chi_B}$.

\begin{lemma}\label{lem:Cperp}
For $z\in C$: (a) $T_zC=\bigcap_A\ker\dd\chi_A(z)$; (b) $(T_zC)^\omega=\operatorname{span}\{X_{\chi_A}(z)\}$; (c) $T_zC\cap(T_zC)^\omega=\{c^AX_{\chi_A}(z):\ c^A\Delta_{AB}(z)=0\ \forall B\}$, of dimension $m-\rank\Delta(z)$.
\end{lemma}
\begin{proof}
(a) is the definition. (b) For $v\in T_zC$, $\omega(X_{\chi_A},v)=\dd\chi_A(v)=0$, so the span lies in $(T_zC)^\omega$; both have dimension $m$ (the $X_{\chi_A}$ are independent because $\omega$ is nondegenerate). (c) By (b) an element of $(T_zC)^\omega$ is $v=c^AX_{\chi_A}$, and $v\in T_zC$ iff $0=\dd\chi_B(v)=c^A\omega(X_{\chi_B},X_{\chi_A})=c^A\Delta_{BA}$.
\end{proof}

\begin{definition}[Dirac's classification]\label{def:class}
The constraint set is \emph{first class} if $\Delta_{AB}\weq0$ (``$\weq$'' denotes equality on $C$), which by Lemma~\ref{lem:Cperp} means $C$ is coisotropic; it is \emph{second class} if $\det\Delta_{AB}\ne0$ on $C$, i.e.\ $C$ is symplectic. In general one can (locally) choose the $\chi_A$ so that $m-\rank\Delta$ of them are first class and $\rank\Delta$ are second class. A function $F$ is first class if $\pb F{\chi_A}\weq0$ for all $A$.
\end{definition}

Thus $\iota_C^*\omega$ has a kernel of dimension $m-\rank\Delta$: it is $0$ for second-class constraints and has dimension $m$ for first-class ones.

%======================================================================
\section{Reduction by constraints}\label{sec:constraints}
%======================================================================

\subsection{Holonomic constraints in Lagrangian form}\label{sec:hololag}

Let $f=(f_1,\dots,f_k):Q\to\R^k$ be smooth with $0$ a regular value. Then $Q_c=f^{-1}(0)$ is an embedded submanifold of dimension $n-k$ (time-independent \emph{holonomic} constraints; the time-dependent, ``rheonomic'' case adds no new ideas). The \emph{virtual displacements} at $q\in Q_c$ are the vectors in $T_qQ_c=\bigcap_a\ker\dd f_a(q)$. The \emph{Lagrange--d'Alembert principle} states that a curve $q(t)\in Q_c$ is a motion of the constrained system if
\begin{equation}\label{eq:dAlembert}
\Big(\frac{\dd}{\dd t}\frac{\partial L}{\partial\dot q^i}-\frac{\partial L}{\partial q^i}\Big)\,\delta q^i=0\qquad\forall\,\delta q\in T_{q(t)}Q_c ,
\end{equation}
that is, the reaction forces do no virtual work (\emph{ideal} constraints). Since \eqref{eq:dAlembert} says that the covector in brackets annihilates $\bigcap_a\ker\dd f_a$, it is a combination of the $\dd f_a$ and \eqref{eq:dAlembert} is equivalent to Lagrange's multiplier rule
\begin{equation}\label{eq:lagmult}
\frac{\dd}{\dd t}\frac{\partial L}{\partial\dot q^i}-\frac{\partial L}{\partial q^i}=\lambda^a\frac{\partial f_a}{\partial q^i},\qquad f(q)=0 .
\end{equation}
The right-hand side is the \emph{constraint force}; the $\lambda^a$ are unknowns determined together with $q(t)$.

\begin{theorem}[Holonomic reduction, Lagrangian form]\label{thm:holored}
Let $L_c=L|_{TQ_c}$. A curve in $Q_c$ satisfies \eqref{eq:dAlembert} if and only if it satisfies the Euler--Lagrange equations of $L_c$ on $Q_c$ (equivalently, is an extremal of the action of $L_c$).
\end{theorem}
\begin{proof}
Choose adapted coordinates $(x^\alpha,y^a)$, $\alpha=1,\dots,n-k$, with $y^a=f_a$ (possible since $0$ is a regular value). Then $Q_c=\{y=0\}$, $TQ_c=\{y=0,\dot y=0\}$ and the virtual displacements are those with $\delta y=0$. Equation~\eqref{eq:dAlembert} is thus the statement that the Euler--Lagrange expressions $E_\alpha(L)=\frac{\dd}{\dd t}\partial L/\partial\dot x^\alpha-\partial L/\partial x^\alpha$ vanish on $\{y=0,\dot y=0\}$. Since restriction to $\{y=0,\dot y=0\}$ commutes with differentiation in $x,\dot x$, $E_\alpha(L)|_{y=0}=E_\alpha(L_c)$. The remaining components give the multipliers, $\lambda_a=E_{y^a}(L)|_{y=0}$, which are determined algebraically once $x(t)$ is known.
\end{proof}

Theorem~\ref{thm:holored} says that on the Lagrangian side holonomic constraint reduction is \emph{trivial}: one restricts the function $L$ to $TQ_c$. (For nonholonomic constraints this fails; Section~\ref{sec:nonhol}.) Constraint forces do no work, so the energy of $L_c$ is conserved.

\subsection{Constrained Hamiltonian dynamics: Dirac's theory}\label{sec:dirac}

Let $(M,\omega)$ be symplectic, $H\in C^\infty(M)$ and $C=\{\chi_A=0\}$ as in Section~\ref{sec:subman}. We write $\iota_C:C\hookrightarrow M$ and look for motions confined to $C$. The following equation is the intrinsic form of constrained dynamics:
\begin{equation}\label{eq:presymp}
X\in\Gamma(TC),\qquad \iota_X(\iota_C^*\omega)=\dd(H|_C).
\end{equation}
Dirac's formulation instead adds multipliers: $X=X_{H+u^A\chi_A}$ on $C$, with $u^A$ chosen so that $X$ is tangent to $C$. These are equivalent.

\begin{proposition}\label{prop:equiv}
A vector field $X$ tangent to $C$ solves \eqref{eq:presymp} if and only if $X=X_H+\lambda^AX_{\chi_A}$ on $C$ for some functions $\lambda^A$. The term $F_C:=\lambda^AX_{\chi_A}$ is the \emph{constraint force}. A solution exists at $z\in C$ if and only if $\dd H(z)$ annihilates $K_z:=T_zC\cap(T_zC)^\omega=\ker\iota_C^*\omega$, and it is then unique modulo $K_z$.
\end{proposition}
\begin{proof}
If $X\in TC$ solves \eqref{eq:presymp} then for every $w\in TC$, $\omega(X,w)=\dd H(w)$, so $\iota_X\omega-\dd H$ annihilates $TC$, hence equals $\lambda_A\,\dd\chi_A$ (using $T_zC=\bigcap\ker\dd\chi_A$) and, by \eqref{eq:XH}, $X=X_H+\lambda^AX_{\chi_A}$ with $\lambda^A=\lambda_A$. Conversely, for $X=X_H+\lambda^AX_{\chi_A}\in TC$ and $w\in TC$: $\omega(X,w)=\dd H(w)+\lambda^A\dd\chi_A(w)=\dd H(w)$. For the last assertion: the map $v\mapsto\iota_v\iota_C^*\omega$ from $T_zC$ to $T_z^*C$ is skew, so its image is the annihilator of its kernel $K_z$; thus $\dd(H|_C)$ lies in the image iff it vanishes on $K_z$, and solutions differ by elements of $K_z$.
\end{proof}

Thus the constrained flow is generated by the \emph{total Hamiltonian} $H_T=H+\lambda^A\chi_A$; the force $F_C$ vanishes at $z$ precisely when $X_H(z)$ is already tangent to $C$.

\begin{theorem}[Dirac bracket; second-class constraints]\label{thm:dirac}
Suppose the constraints are second class, $\det\Delta\neq0$ on $C$ (so $C$ is symplectic and $K=0$). Then \eqref{eq:presymp} has the unique solution
\begin{equation}\label{eq:dirac-sol}
X=X_H+\lambda^AX_{\chi_A},\qquad \lambda^A=-(\Delta^{-1})^{AB}\pb{\chi_B}{H},
\end{equation}
and for every $F\in C^\infty(M)$, $X[F]\weq\pb FH_D$, where the \emph{Dirac bracket} is
\begin{equation}\label{eq:dirac-bracket}
\pb FG_D=\pb FG-\pb F{\chi_A}(\Delta^{-1})^{AB}\pb{\chi_B}G .
\end{equation}
The Dirac bracket restricts to the Poisson bracket of the symplectic manifold $(C,\iota_C^*\omega)$.
\end{theorem}
\begin{proof}
Tangency $X[\chi_B]=0$ gives $\pb{\chi_B}H+\lambda^A\pb{\chi_B}{\chi_A}=0$, i.e.\ $\Delta_{BA}\lambda^A=-\pb{\chi_B}H$, which determines $\lambda$ uniquely. Then $X[F]=\pb FH+\lambda^A\pb F{\chi_A}$ is \eqref{eq:dirac-bracket} with $G=H$. By Proposition~\ref{prop:equiv}, $X$ is the Hamiltonian vector field of $H|_C$ for $\iota_C^*\omega$, so $X[F]|_C=\{F|_C,H|_C\}_{\iota_C^*\omega}$ and the last assertion follows (taking $G$ arbitrary in place of $H$).
\end{proof}

\begin{remark}[First-class constraints and gauge freedom]\label{rem:firstclass}
If $C$ is coisotropic then $K=(TC)^\omega=\operatorname{span}\{X_{\chi_A}\}$ and \eqref{eq:presymp} is solvable iff $\pb H{\chi_A}\weq0$. Then $X_H$ itself is tangent to $C$ and the multipliers $\lambda^A$ are \emph{arbitrary}: the motion is determined only up to the action of the flows $X_{\chi_A}$, which are \emph{gauge transformations} \cite{dirac1964,henneaux1992}. The physical (reduced) phase space is the leaf space $C/K$ (Section~\ref{sec:unified}). If the solvability condition fails, $\pb{H}{\chi_A}$ is a new (``secondary'') constraint and the procedure is iterated: this is the Dirac--Bergmann algorithm \cite{bergmann1949,andersonbergmann1951,dirac1950}, in geometric form the Gotay--Nester--Hinds algorithm \cite{gnh1978,gotaynester1979}.
\end{remark}

\subsection{Holonomic constraints in Hamiltonian form}\label{sec:holoham}

Let $L$ be the mechanical Lagrangian \eqref{eq:mechL} and $f:Q\to\R^k$ as in Section~\ref{sec:hololag}. In $T^*Q$ the correct constraint set is not merely $\{f=0\}$. Put
\begin{equation}\label{eq:chi-phi}
\chi_a:=f_a(q),\qquad \varphi_a:=g^{ij}\,\partial_if_a\,p_j=\dd f_a(p^\sharp),\qquad a=1,\dots,k,
\end{equation}
where $p^\sharp=g\inv p\in TQ$ is the velocity corresponding to $p$, and let $S=\{\chi_a=0,\ \varphi_a=0\}\subset T^*Q$. In words: $q$ lies on $Q_c$ \emph{and the velocity $p^\sharp$ is tangent to $Q_c$}.

\begin{proposition}\label{prop:holoham}
\begin{enumerate}[label=(\roman*)]
\item $S$ is a submanifold of dimension $2(n-k)$, and $\Psi:S\to T^*Q_c$, $\Psi(q,p)=(q,\,p|_{T_qQ_c})$, is a diffeomorphism with $\Psi^*\theta_{Q_c}=\theta|_S$. Hence $\iota_S^*\omega=\Psi^*\omega_{Q_c}$ is symplectic.
\item The Poisson brackets of the constraints are
\begin{equation}\label{eq:holoPB}
\pb{\chi_a}{\chi_b}=0,\qquad\pb{\chi_a}{\varphi_b}=\mathsf G_{ab}:=g^{ij}\partial_if_a\,\partial_jf_b,\qquad \pb{\varphi_a}{\varphi_b}=:B_{ab}.
\end{equation}
The matrix $\mathsf G=(\mathsf G_{ab})$ is symmetric positive definite, so the constraints are second class, with
\begin{equation}\label{eq:holoDelta}
\Delta=\begin{pmatrix}0&\mathsf G\\-\mathsf G&B\end{pmatrix},\qquad
\Delta^{-1}=\begin{pmatrix}\mathsf G\inv B\mathsf G\inv&-\mathsf G\inv\\ \mathsf G\inv&0\end{pmatrix}.
\end{equation}
\item The Legendre map of $L$ maps $\{(q,v):q\in Q_c,\,v\in TQ_c\}$ onto $S$, and the solutions of \eqref{eq:lagmult} (equivalently, of \eqref{eq:dAlembert}) correspond to the solutions of \eqref{eq:presymp} on $S$ with $\chi=(\chi_a,\varphi_a)$. In particular the constrained dynamics is the Hamiltonian dynamics on $T^*Q_c$ of the Legendre transform $H_c$ of $L_c$, and $H_c\circ\Psi=H|_S$.
\end{enumerate}
\end{proposition}
\begin{proof}
(i) The $2k$ functions $\chi_a,\varphi_a$ have independent differentials: $\dd\chi_a=\partial_if_a\dd q^i$ are independent, and $\partial\varphi_a/\partial p_j=g^{ij}\partial_if_a$ are independent since $g$ is nondegenerate. The condition $\varphi_a=0$ says $p^\sharp\in\bigcap\ker\dd f_a=T_qQ_c$. The covector $p|_{T_qQ_c}$ determines $p^\sharp\in T_qQ_c$ because $g|_{T_qQ_c}$ is nondegenerate, so $\Psi$ is bijective with smooth inverse $(q,\alpha)\mapsto(q,g(w,\cdot))$, $w\in T_qQ_c$ the $g$-dual of $\alpha$. For $V\in T_{(q,p)}S$, $\theta(V)=p(T\pi V)$ with $T\pi V\in T_qQ_c$, and $\theta_{Q_c}(T\Psi V)=p|_{T_qQ_c}(T\pi V)$: these agree because $\pi_{Q_c}\circ\Psi=\pi$.
(ii) Direct computation from \eqref{eq:PB}: $\chi$'s depend on $q$ only; $\pb{\chi_a}{\varphi_b}=\partial_kf_a\,\partial\varphi_b/\partial p_k=\partial_kf_a\,g^{kj}\partial_jf_b$. The block form of $\Delta^{-1}$ is verified by multiplication, using $\mathsf G=\mathsf G^T$ and $B=-B^T$.
(iii) The Legendre map is $p=gv$, so $p^\sharp=v$; hence $(q,gv)\in S$ iff $q\in Q_c$ and $v\in T_qQ_c$. Along a solution of \eqref{eq:lagmult}, $\dot q\in TQ_c$, so $(q,p)\in S$. Write \eqref{eq:lagmult} as $\dot q=\partial H/\partial p$, $\dot p_i=-\partial_iH+\lambda^a\partial_if_a$; this is $X_{H-\lambda^a\chi_a}$, since the $\chi_a$ depend on $q$ only. Conversely apply Proposition~\ref{prop:equiv} and Theorem~\ref{thm:dirac} to the $2k$ constraint functions $(\chi_a,\varphi_a)$. By \eqref{eq:holoDelta} the multipliers attached to $\varphi_a$ are $-(\mathsf G\inv)_{ab}\pb{\chi_b}{H}=-(\mathsf G\inv)_{ab}\varphi_b\weq0$, so only the multipliers along $X_{\chi_a}$ survive, which is exactly the multiplier form above. Finally $H|_S(q,p)=\tfrac12g(p^\sharp,p^\sharp)+V=\tfrac12 g_c^{-1}(\alpha,\alpha)+V=H_c(\Psi(q,p))$ with $\alpha=p|_{TQ_c}$.
\end{proof}

\begin{remark}[Naive restriction fails]\label{rem:naiverest}
The set $T^*Q|_{Q_c}=\{\chi=0\}$ has dimension $2n-k$, and $H$ restricted to it still depends on the normal momenta, which are \emph{not} free: they are fixed (to zero, in the metric sense) by $\varphi=0$. The correct reduced phase space has dimension $2n-2k$, and $2k$ constraints are needed. Each holonomic constraint removes two dimensions of phase space: one configuration variable and one momentum.
\end{remark}

By \eqref{eq:dirac-sol} and \eqref{eq:holoDelta} the multipliers are $\lambda^b=(\mathsf G\inv)^{ba}\pb{\varphi_a}{H}$ (up to the sign convention) on $S$: they are \emph{algebraic} functions of the reduced state, and no integration is required to recover the constraint force.

\subsection{Nonholonomic constraints}\label{sec:nonhol}

A constraint is \emph{nonholonomic} if it is not equivalent to a restriction on configurations. The standard case is linear in the velocities,
\begin{equation}\label{eq:nonhol}
\alpha^a_i(q)\,\dot q^i=0,\qquad a=1,\dots,k,
\end{equation}
defining a rank-$(n-k)$ distribution $\mathcal D\subset TQ$; it is holonomic iff $\mathcal D$ is integrable (Frobenius). The Lagrange--d'Alembert principle \eqref{eq:dAlembert} (with $\delta q\in\mathcal D_q$) gives
\begin{equation}\label{eq:nonholeq}
\frac{\dd}{\dd t}\frac{\partial L}{\partial\dot q^i}-\frac{\partial L}{\partial q^i}=\lambda_a\alpha^a_i,\qquad \alpha^a_i\dot q^i=0 .
\end{equation}
For mechanical $L$ these are equations on $N=g^\flat(\mathcal D)\subset T^*Q$, namely $\iota_X\omega=\dd H-\lambda_a\pi_Q^*\alpha^a$ with $\lambda_a$ fixed by tangency to $N$. Three features distinguish the nonholonomic case \cite{akn,bloch2003,cortes2002,bkmm1996}:
\begin{enumerate}[label=(\alph*)]
\item Equations \eqref{eq:nonholeq} are \emph{not} the Euler--Lagrange equations of any restriction of $L$. The variational problem with the constraints imposed on the varied curves (``vakonomic'' mechanics) yields different equations; for holonomic constraints the two approaches agree, for nonholonomic ones they do not, and d'Alembert's principle is the one that agrees with experiment for the usual mechanical constraints \cite{akn}.
\item The dynamics on $N$ is in general not Hamiltonian for any symplectic structure; it is described by an \emph{almost-Poisson} bracket for which the Jacobi identity fails \cite{vanderschaft1994}. The energy is conserved ($\lambda_a\alpha^a_i\dot q^i=0$) but there is in general no invariant symplectic form, and sometimes no invariant measure.
\item Symmetry reduction can still be applied (Section~\ref{sec:nonholsym}), but Noether's theorem holds only in a modified form.
\end{enumerate}
Consequently the presymplectic reduction picture of Section~\ref{sec:unified} does not apply, and nonholonomic reduction is a genuinely different subject.

\subsection{Constraints as limits of stiff potentials}\label{sec:stiff}

Ideal constraints are idealizations of strong forces. Let $H_\kappa=\tfrac12 g^{ij}p_ip_j+V(q)+\tfrac\kappa2\sum_af_a(q)^2$. A classical result \cite{rubin1957} states that for initial data on $S$ the solutions of $H_\kappa$ converge, as $\kappa\to\infty$ and on finite time intervals, to the solution of the constrained system of Proposition~\ref{prop:holoham}, with the constraint force $\lambda^a\partial_if_a$ obtained as a limit of the stiff forces $-\kappa f_a\partial_if_a$. However, the limit is not unique in general: energy stored in the fast transverse oscillations (frequency $\propto\sqrt\kappa$) is an adiabatic invariant, and if it is retained one obtains, in the limit, the constrained motion in an additional effective potential given by the transverse action times the transverse frequency \cite{takens1980,vankampen1984,bornemann1997}. The ideal-constraint model corresponds to vanishing transverse action. In quantum mechanics the analogous limit produces the ``geometric'' da Costa potential \cite{dacosta1981}. Hence constraint reduction embodies a \emph{modelling choice}, not a consequence of the unconstrained dynamics; this is a theme we return to in Section~\ref{sec:compare}.

%======================================================================
\section{Reduction by symmetry}\label{sec:symmetry}
%======================================================================

\subsection{Cyclic coordinates and Routh reduction (abelian case)}\label{sec:routh}

Suppose $Q$ has coordinates $(r^\alpha,x^a)$, $\alpha=1,\dots,n-k$, $a=1,\dots,k$, and $L$ does not depend on the $x^a$; i.e.\ $L$ is invariant under the translation group $G=\R^k$ (or $T^k$) acting on the $x$-coordinates. Take $L$ mechanical, with metric blocks $g_{\alpha\beta}(r)$, $g_{\alpha a}(r)$, $g_{ab}(r)$ and potential $V(r)$. We write $g^{ab}$ for the \emph{inverse of the matrix} $(g_{ab})$ (not a block of $g\inv$), and define
\begin{equation}\label{eq:routhdata}
A^a_\alpha:=g^{ab}g_{b\alpha},\qquad
\mathsf g_{\alpha\beta}:=g_{\alpha\beta}-g_{\alpha a}g^{ab}g_{b\beta},\qquad
V_\mu:=V+\tfrac12g^{ab}\mu_a\mu_b .
\end{equation}
$A^a_\alpha\,\dd r^\alpha$ is the \emph{mechanical connection}, and $\mathsf g$ is positive definite (it is the metric on the shape space $S=Q/G$ for which horizontal vectors are orthogonal to the group orbits).

\begin{proposition}[Routh reduction]\label{prop:routh}
\begin{enumerate}[label=(\alph*)]
\item The momenta $p_a=\partial L/\partial\dot x^a$ are conserved along solutions of \eqref{eq:EL}.
\item Fix $\mu\in\R^k$ and let $\dot x^a(r,\dot r;\mu)$ be the unique solution of $\partial L/\partial\dot x^a=\mu_a$. The \emph{Routhian} is
\begin{equation}\label{eq:Routhian}
R^\mu(r,\dot r):=\big[L-\mu_a\dot x^a\big]_{\dot x=\dot x(r,\dot r;\mu)}
=\tfrac12\mathsf g_{\alpha\beta}\dot r^\alpha\dot r^\beta+\mu_aA^a_\alpha\dot r^\alpha-V_\mu(r).
\end{equation}
A curve $(r(t),x(t))$ with $p_a=\mu_a$ solves \eqref{eq:EL} if and only if $r(t)$ solves the Euler--Lagrange equations of $R^\mu$ on $S$ and $\dot x^a=g^{ab}\mu_b-A^a_\alpha\dot r^\alpha$.
\item In Hamiltonian form, with $\pi_\alpha$ the momentum conjugate to $r^\alpha$ and $\mathsf g^{\alpha\beta}$ the inverse of $\mathsf g_{\alpha\beta}$, the reduced Hamiltonian is
\begin{equation}\label{eq:Hmu}
H_\mu(r,\pi)=H(r,\pi,p=\mu)=\tfrac12\mathsf g^{\alpha\beta}(\pi_\alpha-\mu_aA^a_\alpha)(\pi_\beta-\mu_bA^b_\beta)+V_\mu(r),
\end{equation}
on $(T^*S,\omega_{\rm can})$: a particle on shape space in the ``magnetic field'' $B_\mu=\mu_a\,\dd A^a$ and the amended potential $V_\mu$.
\end{enumerate}
\end{proposition}
\begin{proof}
(a) $\dot p_a=\partial L/\partial x^a=0$. (b) Put $u_a=\mu_a-g_{a\alpha}\dot r^\alpha=g_{ab}\dot x^b$ (this is $\partial L/\partial\dot x^a=\mu_a$ solved for $\dot x$). Then $g_{\alpha a}\dot r^\alpha\dot x^a=(\mu_a-u_a)\dot x^a$ and $\tfrac12g_{ab}\dot x^a\dot x^b=\tfrac12u_a\dot x^a$, so
\[
L-\mu_a\dot x^a=\tfrac12g_{\alpha\beta}\dot r^\alpha\dot r^\beta-\tfrac12g^{ab}u_au_b-V ,
\]
and expanding $u_a$ gives \eqref{eq:Routhian}. Since $\partial L/\partial\dot x^a-\mu_a=0$ at $\dot x=\dot x(r,\dot r;\mu)$, the chain rule gives
$\partial R^\mu/\partial\dot r^\alpha=\partial L/\partial\dot r^\alpha$ and $\partial R^\mu/\partial r^\alpha=\partial L/\partial r^\alpha$ (the derivatives of $\dot x(r,\dot r;\mu)$ drop out). Hence the $r$-components of \eqref{eq:EL} are the Euler--Lagrange equations of $R^\mu$; the $x$-components are $\dot p_a=0$ and the equation for $\dot x$ above. (c) Inverting the block metric gives $g^{-1}$ with blocks $\mathsf g^{\alpha\beta}$, $-\mathsf g^{\alpha\beta}A^b_\beta$ and $g^{ab}+A^a_\alpha\mathsf g^{\alpha\beta}A^b_\beta$, hence $H=\tfrac12\mathsf g^{\alpha\beta}(\pi_\alpha-p_aA^a_\alpha)(\pi_\beta-p_bA^b_\beta)+\tfrac12g^{ab}p_ap_b+V$; set $p_a=\mu_a$. One checks that $R^\mu=\pi_\alpha\dot r^\alpha-H_\mu$ with $\dot r^\alpha=\partial H_\mu/\partial\pi_\alpha$, i.e.\ $R^\mu$ is the Legendre transform of $H_\mu$.
\end{proof}

\begin{remark}[Reconstruction is a quadrature]\label{rem:quad}
Once $r(t)$ is known, $x^a(t)=x^a(0)+\int_0^t\big(g^{ab}\mu_b-A^a_\alpha\dot r^\alpha\big)\dd t$. The group variables are recovered by \emph{integration}, and the result splits into a dynamic part (from $g^{ab}\mu_b$) and a geometric part (the holonomy of the connection $A$ along the loop traced in shape space) \cite{mmr1990,montgomery1991}.
\end{remark}

\begin{remark}[The Routh pitfall]\label{rem:pitfall1}
One cannot reduce by substituting $\dot x(r,\dot r;\mu)$ into $L$ and then varying: $L|_{\dot x=\dot x(\cdot;\mu)}=R^\mu+\mu_a\dot x^a(r,\dot r;\mu)$, and the second term has $r$- and $\dot r$-dependent derivatives that spoil the equations. The correct procedure subtracts $\mu_a\dot x^a$. (Example~\ref{ex:kepler} shows the sign error this produces.) By contrast, in the Hamiltonian formulation one simply sets $p_a=\mu_a$ in $H$.
\end{remark}

\subsection{The Marsden--Weinstein--Meyer theorem}\label{sec:MWM}

\begin{lemma}\label{lem:kernel}
Let $G$ act freely and symplectically on $(M,\omega)$ with equivariant momentum map $J$. Then $J$ is a submersion, so $J^{-1}(\mu)$ is a closed embedded submanifold of codimension $\dim G$, and for $z\in J^{-1}(\mu)$:
\begin{enumerate}[label=(\alph*)]
\item $T_zJ^{-1}(\mu)=\ker T_zJ=\big(T_z(G\cdot z)\big)^\omega$;
\item $\ker\big(\iota_\mu^*\omega\big)_z=T_zJ^{-1}(\mu)\cap T_z(G\cdot z)=T_z(G_\mu\cdot z)$,
\end{enumerate}
where $\iota_\mu:J^{-1}(\mu)\hookrightarrow M$.
\end{lemma}
\begin{proof}
By \eqref{eq:momap}, $\pair{T_zJ\,v}{\xi}=\omega(\xi_M(z),v)$ for all $v\in T_zM$, $\xi\in\frakg$. A $\xi$ annihilating the image of $T_zJ$ therefore satisfies $\xi_M(z)\in(T_zM)^\omega=0$, so $\xi=0$ for a free action: $T_zJ$ is onto, $J^{-1}(\mu)$ is a submanifold and $T_zJ^{-1}(\mu)=\ker T_zJ$. The displayed identity also shows $v\in\ker T_zJ$ iff $\omega(\xi_M,v)=0$ for all $\xi$, i.e.\ $v\in(T_z(G\cdot z))^\omega$; this is (a). For (b), the kernel of the restriction of $\omega$ to $W=\ker T_zJ$ is $W\cap W^\omega=\ker T_zJ\cap T_z(G\cdot z)$, since $W^\omega=(T_z(G\cdot z))^{\omega\omega}=T_z(G\cdot z)$. Finally, $\xi_M(z)\in\ker T_zJ$ iff $\ad^*_\xi\mu=0$ by \eqref{eq:TJ}, i.e.\ iff $\xi\in\frakg_\mu$.
\end{proof}

\begin{theorem}[Marsden--Weinstein--Meyer \cite{mw1974,meyer1973}]\label{thm:MWM}
Let $G$ act freely and properly on $(M,\omega)$ by a Hamiltonian action with equivariant momentum map $J$, and let $\mu\in\frakg^*$. Then:
\begin{enumerate}[label=(\alph*)]
\item $G_\mu$ acts freely and properly on $J^{-1}(\mu)$, and $M_\mu:=J^{-1}(\mu)/G_\mu$ is a smooth manifold of dimension
\begin{equation}\label{eq:dimMmu}
\dim M_\mu=\dim M-\dim G-\dim G_\mu .
\end{equation}
\item There is a unique symplectic form $\omega_\mu$ on $M_\mu$ with $\pi_\mu^*\omega_\mu=\iota_\mu^*\omega$.
\item If $H$ is $G$-invariant, the flow $\varphi_t$ of $X_H$ preserves $J^{-1}(\mu)$, commutes with $G$, and therefore induces a flow $\varphi^\mu_t$ on $M_\mu$ with $\pi_\mu\circ\varphi_t=\varphi^\mu_t\circ\pi_\mu$. This is the Hamiltonian flow of $H_\mu$, where $H_\mu\circ\pi_\mu=H\circ\iota_\mu$.
\end{enumerate}
\end{theorem}
\begin{proof}
(a) $G_\mu$ is a closed subgroup of $G$, hence acts freely and properly on $M$, and it preserves the closed submanifold $J^{-1}(\mu)$ because $J(g\cdot z)=\Ad^*_g\mu=\mu$ for $g\in G_\mu$. The dimension count is $\dim M-\dim G-\dim G_\mu$ by Lemma~\ref{lem:kernel}.
(b) By Lemma~\ref{lem:kernel}(b), $\ker\iota_\mu^*\omega=\ker T\pi_\mu$, the tangent space of the $G_\mu$-orbit. Hence $\iota_\mu^*\omega$ is basic with respect to $\pi_\mu$: it is invariant under $G_\mu$ and annihilated by the vertical vectors. It therefore descends to a form $\omega_\mu$ with $\pi_\mu^*\omega_\mu=\iota_\mu^*\omega$. It is closed since $\pi_\mu^*\dd\omega_\mu=\dd\iota_\mu^*\omega=0$ and $\pi_\mu$ is a surjective submersion; it is nondegenerate because the kernel of its pullback is exactly the kernel of $T\pi_\mu$.
(c) By Proposition~\ref{prop:noether} $J$ is constant along $X_H$, so $X_H$ is tangent to $J^{-1}(\mu)$; invariance of $H$ and $\omega$ makes $\varphi_t$ commute with $G$, in particular with $G_\mu$, so it maps $G_\mu$-orbits to $G_\mu$-orbits and descends. $H\circ\iota_\mu$ is $G_\mu$-invariant, so $H_\mu$ is well defined. For $v\in T_zJ^{-1}(\mu)$: $\omega(X_H,v)=\dd H(v)=\dd(H_\mu\circ\pi_\mu)(v)=\dd H_\mu(T\pi_\mu v)$, whence $\iota_{T\pi_\mu X_H}\omega_\mu=\dd H_\mu$.
\end{proof}

\begin{corollary}[Abelian case]\label{cor:abelian}
If $G$ is abelian then $G_\mu=G$, $J^{-1}(\mu)$ is coisotropic with characteristic leaves equal to the $G$-orbits, and $\dim M_\mu=\dim M-2\dim G$.
\end{corollary}
\begin{proof}
For abelian $G$, $\ad^*\equiv0$ so $G_\mu=G$, and Lemma~\ref{lem:kernel}(b) gives $\ker\iota_\mu^*\omega=T(G\cdot z)=(TJ^{-1}(\mu))^\omega$.
\end{proof}

\paragraph{Reconstruction.} Given a trajectory $\bar z(t)$ of $X_{H_\mu}$ and a point $z_0\in\pi_\mu^{-1}(\bar z(0))$ there is a unique trajectory of $X_H$ through $z_0$, and it projects to $\bar z(t)$. Choosing a connection on the principal $G_\mu$-bundle $J^{-1}(\mu)\to M_\mu$ one writes it as a horizontal lift composed with a motion $g(t)$ in $G_\mu$ obtained by integrating an equation driven by $\bar z(t)$ (the reconstruction equation); the total phase is the sum of a dynamic and a geometric (holonomy) part \cite{mmr1990,montgomery1991}. Thus information lost in symmetry reduction---the group variables---is recovered by quadrature.

\subsection{Cotangent bundles, Lie--Poisson and Euler--Poincar\'e reduction}\label{sec:cotred}

For $M=T^*Q$ with the cotangent lifted action the reduced space has a concrete description. For abelian $G$ acting freely, Proposition~\ref{prop:routh}(c) is the statement that $(T^*Q)_\mu\cong T^*(Q/G)$ with the canonical form in the variables $(r,\pi)$, but the Hamiltonian acquires the magnetic term $\mu_aA^a_\alpha$. In terms of the \emph{kinetic} momenta $\pi_\alpha-\mu_aA^a_\alpha$ one equivalently keeps $H$ in its standard form and replaces $\omega_{\rm can}$ by $\omega_{\rm can}-B_\mu$ (cotangent bundle reduction with a ``magnetic term'' \cite{smale1970,satzer1977,kummer1981,marsden1992}); in general the reduced space $(T^*Q)_\mu$ is a symplectic fibre bundle over $T^*(Q/G)$ with fibre the coadjoint orbit $\mathcal O_\mu$ \cite{mr1999,marsdenperlmutter2000}. In particular, \emph{the reduced symplectic form depends on $\mu$}: the reduced spaces for different $\mu$ are different symplectic manifolds, glued together into a Poisson manifold $M/G$.

For $Q=G$ with $G$ acting on itself by left translations, $(T^*G)/G\cong\frakg^*$ with the Lie--Poisson bracket $\pb FG_\pm(\mu)=\pm\pair\mu{[\delta F/\delta\mu,\delta G/\delta\mu]}$, whose symplectic leaves are the coadjoint orbits $\mathcal O_\mu=(T^*G)_\mu$ \cite{am1978,mr1999}. The Lagrangian counterpart is Euler--Poincar\'e reduction: for a left-invariant $L:TG\to\R$ with $\ell=L|_{\frakg}$ the equations \eqref{eq:EL} reduce to
\begin{equation}\label{eq:EP}
\frac{\dd}{\dd t}\frac{\delta\ell}{\delta\xi}=\ad^*_\xi\frac{\delta\ell}{\delta\xi},
\end{equation}
found by Poincar\'e \cite{poincare1901}, see \cite{hmr1998}. For the free rigid body $\frakg=\mathfrak{so}(3)\cong\R^3$, $\ell=\frac12\Omega\cdot I\Omega$, $\Pi=I\Omega$, and \eqref{eq:EP} is Euler's equation $\dot\Pi=\Pi\times\Omega$. The general Lagrangian theory on $TQ/G$ (Lagrange--Poincar\'e equations) and its relation to Routh reduction are in \cite{marsdenscheurle1993,mrs2000,cmr2001}; reduction in stages is treated in \cite{mmopr2007}.

\subsection{Singular reduction}\label{sec:singular}

If $G$ does not act freely, or $\mu$ is not a regular value of $J$, then $J^{-1}(\mu)/G_\mu$ need not be a manifold. For proper actions it is a \emph{stratified symplectic space}: a disjoint union of symplectic manifolds (one for each orbit type) fitted together in a controlled way \cite{sjamaarlerman1991,op2004}. The phenomenon is familiar already in spherical coordinates, where the polar coordinate singularity reflects the non-free action of rotations about the vertical axis at the poles of the sphere (Example~\ref{ex:sphpend}).

%======================================================================
\section{A common framework: presymplectic reduction}\label{sec:unified}
%======================================================================

\subsection{The theorem}

\begin{theorem}[Presymplectic reduction \cite{gnh1978,gotay1982,marle1995,mw1974}]\label{thm:presymp}
Let $(M,\omega)$ be symplectic and $\iota:C\hookrightarrow M$ an embedded submanifold such that $\iota^*\omega$ has constant rank, and put $K=\ker\iota^*\omega=TC\cap(TC)^\omega$.
\begin{enumerate}[label=(\alph*)]
\item $K$ is an involutive distribution on $C$.
\item If the leaf space $\bar C=C/K$ is a manifold and $\pi:C\to\bar C$ a submersion, there is a unique symplectic form $\bar\omega$ with $\pi^*\bar\omega=\iota^*\omega$. If $C=\{\chi_A=0\}$ with $m$ constraints, $m-\rank\Delta$ of them first class and $\rank\Delta$ second class (Definition~\ref{def:class}), then
\begin{equation}\label{eq:dimred}
\dim\bar C=\dim M-2(m-\rank\Delta)-\rank\Delta .
\end{equation}
\item Let $H\in C^\infty(M)$ with $\dd H|_K=0$. Then \eqref{eq:presymp} has solutions; the flow of any solution $X$ preserves $\iota^*\omega$, maps leaves to leaves, and projects to the Hamiltonian flow of $\bar H$ on $(\bar C,\bar\omega)$, where $\bar H\circ\pi=H|_C$.
\end{enumerate}
\end{theorem}
\begin{proof}
(a) Let $X,Y\in\Gamma(K)$ and $Z\in\Gamma(TC)$. Since $\dd\iota^*\omega=0$, the usual formula for the exterior derivative of a two-form gives $0=\dd\iota^*\omega(X,Y,Z)=-\omega([X,Y],Z)$, all other terms containing $\omega(X,\cdot)$ or $\omega(Y,\cdot)$ and hence vanishing. Thus $[X,Y]\in K$.
(b) For $X\in\Gamma(K)$, $\mathcal L_X\iota^*\omega=\dd\iota_X\iota^*\omega+\iota_X\dd\iota^*\omega=0$, so $\iota^*\omega$ is basic with respect to $\pi$ and descends to a closed form $\bar\omega$ whose pullback has kernel exactly $\ker T\pi$, hence $\bar\omega$ is nondegenerate. For the dimension: $\dim C=\dim M-m$ and $\dim K=m-\rank\Delta$ by Lemma~\ref{lem:Cperp}(c).
(c) Existence by Proposition~\ref{prop:equiv}. If $X$ solves \eqref{eq:presymp} then $\mathcal L_X\iota^*\omega=\dd(\dd (H|_C))=0$, so the flow preserves $\iota^*\omega$ and therefore its kernel $K$; it maps leaves to leaves and descends. Since $\dd H|_K=0$ (and leaves are connected), $H|_C=\bar H\circ\pi$, and $\pi^*(\iota_{T\pi X}\bar\omega)=\iota_X\iota^*\omega=\dd(H|_C)=\pi^*\dd\bar H$.
\end{proof}

Equation \eqref{eq:dimred} is Dirac's counting rule: each first-class constraint removes two dimensions of phase space (the constraint surface and the gauge orbit), each second-class constraint removes one.

\subsection{Constraint reduction as presymplectic reduction}

\emph{Holonomic constraints.} By Proposition~\ref{prop:holoham}, $C=S\cong T^*Q_c$ is symplectic, $K=0$, $\bar C=S$ and \eqref{eq:dimred} gives $2n-2k$. The reduction is a pure \emph{restriction}: $\bar\omega=\iota_S^*\omega$.
\emph{Gauge (first-class) constraints.} If $C$ is coisotropic, $K=(TC)^\omega$ and $\bar C=C/K$ is obtained by restriction \emph{and} quotient (Remark~\ref{rem:firstclass}); the leaves of $K$ are the gauge orbits.

\subsection{Symmetry reduction as presymplectic reduction}

\begin{proposition}\label{prop:symclass}
Let $G$ act as in Theorem~\ref{thm:MWM}, put $C=J^{-1}(\mu)=\{\chi_a=0\}$ with $\chi_a=J_a-\mu_a$, $a=1,\dots,\dim G$. Then
\begin{equation}\label{eq:DeltaJ}
\Delta_{ab}=\pb{J_a}{J_b}=c^c_{ab}J_c\ \weq\ c^c_{ab}\mu_c=\pair\mu{[e_a,e_b]},
\end{equation}
a constant antisymmetric matrix with kernel $\frakg_\mu$ and rank $\dim\mathcal O_\mu=\dim G-\dim G_\mu$ (it is the Kirillov--Kostant--Souriau form of $\mathcal O_\mu$ at $\mu$). Hence, with $e_1,\dots,e_s$ a basis of $\frakg_\mu$:
\begin{itemize}
\item $J_\xi-\pair\mu\xi$ for $\xi\in\frakg_\mu$ are $s=\dim G_\mu$ \emph{first-class} constraints;
\item the remaining $\dim\mathcal O_\mu$ components are \emph{second class}, and are removed by the Dirac bracket;
\item the kernel of $\iota_\mu^*\omega$ is $\{\xi_M:\xi\in\frakg_\mu\}=T(G_\mu\cdot z)$, and \eqref{eq:dimred} reproduces \eqref{eq:dimMmu}: $2\dim G_\mu+\dim\mathcal O_\mu=\dim G+\dim G_\mu$.
\end{itemize}
In particular, for abelian $G$ all $\dim G$ constraints are first class.
\end{proposition}
\begin{proof}
The first identity is \eqref{eq:TJ}; the kernel statement is \eqref{eq:gmu}. A vector $\xi\in\frakg_\mu$ gives $\pair\mu{[\xi,\eta]}=0$ for all $\eta$, so the corresponding constraint has vanishing bracket with all others on $C$: it is first class. On a complement of $\frakg_\mu$ the matrix is nondegenerate. By Lemma~\ref{lem:Cperp}(c) the kernel of $\iota_\mu^*\omega$ consists of $c^aX_{\chi_a}=(c^ae_a)_M$ with $c^a\Delta_{ab}=0$, i.e.\ $c^ae_a\in\frakg_\mu$. This is Lemma~\ref{lem:kernel}(b) again.
\end{proof}

Hence the Marsden--Weinstein--Meyer theorem is \emph{Dirac's second-class/first-class reduction applied to the constraints $J=\mu$}, followed by division by the gauge flows of the first-class part. The mathematical machine is the same as for ``pure'' constraints; what differs is the origin of the constraint surface and the nature of its dynamics, which we now make precise.

\subsection{Same equation, different solutions}

\begin{theorem}\label{thm:samesolve}
Let $C=\{\chi=0\}$ and let $X$ solve \eqref{eq:presymp}; thus $X=X_H+F_C$ with $F_C=\lambda^AX_{\chi_A}$ (Proposition~\ref{prop:equiv}).
\begin{enumerate}[label=(\alph*)]
\item $F_C$ can be taken to vanish at $z\in C$ if and only if $X_H(z)\in T_zC$, i.e.\ iff $\pb{\chi_A}{H}(z)=0$ for all $A$.
\item \emph{Symmetry reduction:} for $C=J^{-1}(\mu)$ and $H$ $G$-invariant, $\pb{J_a-\mu_a}{H}=0$ \emph{identically}. Hence $X=X_H$ is a solution (zero constraint force), and all solutions differ from it by elements of $K=T(G_\mu\cdot z)$, which project to the same reduced trajectory.
\item \emph{Second-class constraint reduction:} the solution is unique, $F_C=\lambda^AX_{\chi_A}$ with $\lambda^A$ given by \eqref{eq:dirac-sol}, and $F_C\not\equiv0$ in general: it vanishes along $C$ exactly when $C$ is invariant under the flow of $X_H$.
\end{enumerate}
Consequently, the trajectories of a symmetry-reduced system are projections of trajectories of the \emph{unreduced} system; the trajectories of a constrained system are, in general, not trajectories of the unconstrained one.
\end{theorem}
\begin{proof}
(a) $X_H(z)\in T_zC$ iff $\dd\chi_A(X_H)=\pb{\chi_A}{H}=0$; if so $X=X_H$ solves \eqref{eq:presymp}. (b) is Proposition~\ref{prop:noether}. (c) is Theorem~\ref{thm:dirac}: $F_C=0$ iff $\pb{\chi_B}{H}\weq0$ for all $B$ iff $X_H$ is tangent to $C$, iff $C$ is invariant.
\end{proof}

For example, for a free particle in $\R^3$ constrained to a sphere, $X_H$ is a straight line leaving the sphere and $F_C\ne0$; constrained to a plane through the origin, $F_C=0$. For the Kepler problem $F_C=0$ on every $J^{-1}(\mu)$.

\subsection{Reduction commutes with second-class constraints}\label{sec:commute}

When the constraint set and the symmetry are compatible, the two reductions can be performed in either order.

\begin{proposition}\label{prop:commute}
Let $G$ act on $(M,\omega)$ as in Theorem~\ref{thm:MWM} and let $S\subset M$ be a $G$-invariant symplectic submanifold on which $\mu$ is a regular value of $J|_S$ and $G_\mu$ acts freely and properly. Then $G$ acts on $(S,\omega|_S)$ in a Hamiltonian way with momentum map $J|_S$, and the reduced space $S_\mu=(S\cap J^{-1}(\mu))/G_\mu$ is (the image of) a symplectic submanifold of $M_\mu$, whose symplectic form is the restriction of $\omega_\mu$.
\end{proposition}
\begin{proof}
$\xi_M$ is tangent to $S$ because $S$ is invariant, and $\iota_{\xi_S}\omega|_S=\dd(J_\xi|_S)$, so $J|_S$ is an equivariant momentum map. $S\cap J^{-1}(\mu)=(J|_S)^{-1}(\mu)$ is a $G_\mu$-invariant submanifold of $J^{-1}(\mu)$, and the induced map $S_\mu\to M_\mu$ is injective because $G_\mu$-orbits through points of $S\cap J^{-1}(\mu)$ stay in $S$. Pulling back to $S\cap J^{-1}(\mu)$, $\omega_\mu|_{S_\mu}$ satisfies $\pi^*(\omega_\mu|_{S_\mu})=\omega|_{S\cap J^{-1}(\mu)}=\pi^*\omega_{S,\mu}$; so $\omega_\mu|_{S_\mu}=\omega_{S,\mu}$, which is nondegenerate by Theorem~\ref{thm:MWM}(b) applied to $S$.
\end{proof}

This is the situation of the spherical pendulum and the rigid body (Examples~\ref{ex:sphpend} and~\ref{ex:rigid}).

\subsection{Poisson-algebraic remark}

In terms of observables, the three kinds of reduction act differently on the Poisson algebra $C^\infty(M)$. Let $I_C$ be the ideal of functions vanishing on $C$. For a \emph{coisotropic} $C$ the reduced algebra is $C^\infty(\bar C)\cong\{F:\pb F{I_C}\subset I_C\}/I_C$: first-class functions modulo constraints (Dirac observables). For a \emph{second-class} $C$, $C^\infty(C)=C^\infty(M)/I_C$ carries the Dirac bracket \eqref{eq:dirac-bracket} and every function survives. For symmetry reduction, $C^\infty(M_\mu)\cong C^\infty(J^{-1}(\mu))^{G_\mu}$: $G_\mu$-invariant functions restricted to the level set. A common generalization is Poisson reduction of Marsden and Ratiu \cite{mr1986}, which gives conditions on a submanifold $C$ and a distribution $F\subset TC$ under which $C/F$ inherits a Poisson structure and which contains Dirac's bracket and Marsden--Weinstein reduction as special cases.

%======================================================================
\section{Worked examples}\label{sec:examples}
%======================================================================

\begin{example}[Plane pendulum: constraint only]\label{ex:pendulum}
Take $Q=\R^2\setminus\{0\}$ with polar coordinates $(r,\theta)$, $\theta$ measured from the downward vertical, and $L=\tfrac m2(\dot r^2+r^2\dot\theta^2)+mgr\cos\theta$. On $T^*Q$ (dimension $4$),
\[
H=\frac1{2m}\Big(p_r^2+\frac{p_\theta^2}{r^2}\Big)-mgr\cos\theta .
\]
The constraint $f=r-\ell$ gives, in \eqref{eq:chi-phi}, $\chi_1=r-\ell$ and $\chi_2=p_r$ (the latter is $m\varphi$). Then $\pb{\chi_1}{\chi_2}=1$, so $\Delta^{-1}=\begin{psmallmatrix}0&-1\\1&0\end{psmallmatrix}$. Since $\chi_1,\chi_2$ Poisson-commute with $\theta,p_\theta$, the Dirac bracket \eqref{eq:dirac-bracket} of $(\theta,p_\theta)$ is canonical and the reduced Hamiltonian on $S\cong T^*S^1$ is
\[
H_S=\frac{p_\theta^2}{2m\ell^2}-mg\ell\cos\theta,\qquad\text{so}\qquad \ddot\theta=-\frac g\ell\sin\theta .
\]
\emph{Constraint force.} From \eqref{eq:dirac-sol}: $\{\chi_2,H\}=-\partial H/\partial r=p_\theta^2/(mr^3)+mg\cos\theta$, which on $S$ equals the tension
\[
T=m\ell\dot\theta^2+mg\cos\theta .
\]
Then $\lambda^1=T$, $\lambda^2\weq0$, and $F_C=T\,X_{r-\ell}=-T\,\partial/\partial p_r$: it cancels exactly the radial acceleration the free dynamics would produce, $\dot p_r|_{X_H}=+T$. Note (i) $F_C\ne0$, so the reduced motion is not a motion of the unconstrained particle; (ii) $F_C$ is recovered \emph{algebraically} from the reduced state $(\theta,\dot\theta)$; (iii) the single holonomic constraint ($k=1$) removed $2k=2$ phase-space dimensions, $4\to2$.
\end{example}

\begin{example}[Kepler problem: symmetry only]\label{ex:kepler}
Take $L=\tfrac m2(\dot r^2+r^2\dot\theta^2)+k/r$ on $T^*(\R^2\setminus0)$. The rotation group $G=SO(2)$ acts freely by $\theta\mapsto\theta+\text{const}$ with momentum map $J=p_\theta=mr^2\dot\theta=:\mu$, and $G$ is abelian, so Corollary~\ref{cor:abelian} gives $M_\mu=T^*\R_+$, of dimension $4-2=2$. By Proposition~\ref{prop:routh} the Routhian is
\[
R^\mu=\tfrac m2\dot r^2-V_{\rm eff}(r),\qquad V_{\rm eff}(r)=-\frac kr+\frac{\mu^2}{2mr^2},
\]
equivalently $H_\mu=p_r^2/2m+V_{\rm eff}(r)$ (here $\mathsf g=m$, $A=0$ because the metric is diagonal). The centrifugal term arises from the conservation law, \emph{not} from a constraint force: $F_C\equiv0$ in the sense of Theorem~\ref{thm:samesolve}(b), and every reduced trajectory is the projection of a genuine Kepler trajectory. Reconstruction is a quadrature, $\theta(t)=\theta_0+\int_0^t\mu\,\dd t'/(mr(t')^2)$.

\emph{The pitfall of Remark~\ref{rem:pitfall1}.} Substituting $\dot\theta=\mu/(mr^2)$ into $L$ gives $\tfrac m2\dot r^2+\mu^2/(2mr^2)+k/r$, whose Euler--Lagrange equation is $m\ddot r=-\mu^2/(mr^3)-k/r^2$: the centrifugal term has the \emph{wrong sign}. The correct equation, from $R^\mu$, is $m\ddot r=+\mu^2/(mr^3)-k/r^2$.
\end{example}

\begin{example}[Spherical pendulum: constraint and symmetry]\label{ex:sphpend}
Start with a particle in $T^*\R^3$ (dimension $6$) with $H=|p|^2/2m+mgz$. The constraint $f=\tfrac12(|x|^2-\ell^2)$ gives, by \eqref{eq:chi-phi}, $S=\{|x|=\ell,\ x\cdot p=0\}\cong T^*S^2$ of dimension $4$ (two second-class constraints). In spherical coordinates (with $\vartheta$ measured from the downward vertical)
\[
H_S=\frac{p_\vartheta^2}{2m\ell^2}+\frac{p_\varphi^2}{2m\ell^2\sin^2\vartheta}-mg\ell\cos\vartheta .
\]
The group $G=SO(2)$ of rotations about the vertical axis leaves $S$ and $H$ invariant, with $J=p_\varphi=xp_y-yp_x$; by Proposition~\ref{prop:commute} it does not matter whether one imposes the constraint first or reduces first. Reducing $T^*S^2$ at $J=\mu\ne0$ gives $M_\mu=T^*(0,\pi)$ of dimension $4-2=2$ with
\[
H_\mu(\vartheta,p_\vartheta)=\frac{p_\vartheta^2}{2m\ell^2}+V_\mu(\vartheta),\qquad V_\mu(\vartheta)=\frac{\mu^2}{2m\ell^2\sin^2\vartheta}-mg\ell\cos\vartheta,
\]
and $\vartheta$ oscillates between two turning points since $V_\mu\to\infty$ at the poles. The azimuth is reconstructed from $\dot\varphi=\mu/(m\ell^2\sin^2\vartheta)$.

The dimension count is $6\xrightarrow{\ 2\text{ constraints}\ }4\xrightarrow{\ 1\text{ symmetry}\ }2$, with the constraint removing one dimension per second-class function and the symmetry two per abelian generator. The constraint force (tension) is nonzero, while the centrifugal-like term $\mu^2/(2m\ell^2\sin^2\vartheta)$ is not a force of constraint. For $\mu=0$ the isotropy at the poles is non-trivial and one is in the singular regime of Section~\ref{sec:singular}.
\end{example}

\begin{example}[Rigid body: constraint then symmetry]\label{ex:rigid}
Take $N\ge3$ non-collinear particles with all mutual distances fixed. The $3N-6$ independent holonomic constraints reduce $T^*\R^{3N}$ (dimension $6N$) to $T^*SE(3)$ (dimension $12$). Passing to the centre-of-mass frame is an abelian symmetry reduction by translations, giving $T^*SO(3)$ of dimension $6$. For free motion $H=\tfrac12\Pi\cdot I^{-1}\Pi$ is invariant under $G=SO(3)$ acting by left translation (rotations of space), with momentum map the spatial angular momentum $\mu$. Marsden--Weinstein reduction at $\mu\ne0$ has $G_\mu\cong SO(2)$ (rotations about $\mu$), and
\[
\dim M_\mu=6-3-1=2,\qquad M_\mu\cong\mathcal O_\mu=\{\Pi\in\R^3:|\Pi|=|\mu|\},
\]
the sphere in body-angular-momentum space with the area form as symplectic form. By \eqref{eq:EP} the reduced dynamics is Euler's equation $\dot\Pi=\Pi\times I^{-1}\Pi$, whose trajectories are the intersections of the sphere with the energy ellipsoids. Here Proposition~\ref{prop:symclass} reads: $\dim\mathcal O_\mu=2$ second-class and $\dim G_\mu=1$ first-class constraints among the three components of $J-\mu$, in agreement with $2\cdot1+2=3+1$.
\end{example}

%======================================================================
\section{Comparison}\label{sec:compare}
%======================================================================

We now compare the two procedures along the axes raised in (Q1)--(Q3). A summary is given in Table~\ref{tab:summary}.

\subsection{Origin and logical status}

\emph{Constraint reduction is an input; symmetry reduction is a consequence.} A constraint $C$ is part of the \emph{specification} of the model. It is not a property of $L$ or $H$ and cannot be derived from them; it represents an idealization of interactions (rigid links, rails, strong binding forces) whose detailed mechanism has been suppressed. Different realizations of ``the same'' constraint can yield different limiting dynamics (Section~\ref{sec:stiff}). A symmetry, in contrast, is a \emph{property} of the dynamical data $(L\text{ or }H)$. The conserved momenta follow from Noether's theorem (Proposition~\ref{prop:noether}), every value $\mu$ of the momentum map occurs in the unreduced system, and fixing $\mu$ is a choice of initial data (a ``sector''), not a modification of the model.

\subsection{Effect on the dynamics}

By Theorem~\ref{thm:samesolve}, both procedures solve the same equation $\iota_X\iota_C^*\omega=\dd(H|_C)$, but with opposite outcomes for the constraint force $F_C$:
\begin{itemize}
\item \emph{Symmetry:} $C=J^{-1}(\mu)$ is \emph{dynamically invariant} (the unconstrained flow already stays on it), so $F_C=0$ and reduced trajectories are images of genuine trajectories.
\item \emph{Constraint:} $C$ is \emph{imposed}; generically $X_H$ leaves $C$, and $F_C\neq0$ must be added to keep the motion on $C$. The constrained trajectories are not trajectories of the unconstrained system.
\end{itemize}
A corollary concerns what is lost and how it is recovered. In symmetry reduction the lost data (group variables) are recovered from the reduced solution by \emph{quadrature}, with a geometric phase (Remark~\ref{rem:quad}). In constraint reduction the lost data (normal coordinates and momenta) are fixed by hypothesis and are \emph{not} recoverable from the dynamics; what is recovered is the reaction force, \emph{algebraically} (Example~\ref{ex:pendulum}).

\subsection{Geometry: restriction versus restriction and quotient}

Both are instances of Theorem~\ref{thm:presymp}. The distinction is the kernel $K$ of $\iota_C^*\omega$:
\begin{center}
\begin{tabular}{@{}lll@{}}
\toprule
 & $C$ & $K=\ker\iota_C^*\omega$\\
\midrule
holonomic constraints (second class) & $S\cong T^*Q_c$ & $0$\\
gauge constraints (first class) & coisotropic & $(TC)^\omega$\\
symmetry, $G$ abelian & $J^{-1}(\mu)$ & $T(G\cdot z)$\\
symmetry, $G$ non-abelian & $J^{-1}(\mu)$ & $T(G_\mu\cdot z)$\\
\bottomrule
\end{tabular}
\end{center}
Constraint reduction (second class) is therefore a pure \emph{restriction}; symmetry reduction is \emph{restriction followed by a quotient}. The dimension counts agree because of \eqref{eq:dimred}: a holonomic constraint removes two phase-space dimensions, a one-parameter abelian symmetry removes two (fix $p$, forget $x$), and in general $\dim G+\dim G_\mu$ for a non-abelian symmetry.

\subsection{Parameters and the reduced symplectic structure}

Constraint reduction produces a single reduced space with the canonical form of $T^*Q_c$. Symmetry reduction produces a \emph{family} $\{M_\mu\}$ indexed by $\mu\in\frakg^*$, and the symplectic form of $M_\mu$ depends on $\mu$ through the magnetic term $B_\mu$ (Section~\ref{sec:cotred}); the union of the $M_\mu$ is the Poisson manifold $M/G$. The family structure is absent for holonomic constraints (the constraint values are fixed numbers). The ``centrifugal'' and ``gyroscopic'' terms in reduced Hamiltonians---$\mu^2/2mr^2$ in Example~\ref{ex:kepler}, the term $\mu_aA^a_\alpha\dot r^\alpha$ in \eqref{eq:Routhian}---are artifacts of conservation laws, not of reaction forces.

\subsection{Dependence on the formalism: dual pitfalls}

A striking asymmetry emerges between the Lagrangian and Hamiltonian formulations.

\begin{itemize}
\item \emph{Lagrangian picture.} A holonomic constraint is a condition on \emph{points} of $Q$. Reduction is trivial: restrict $L$ to $TQ_c$ (Theorem~\ref{thm:holored}). A conserved momentum is a condition on \emph{velocities}; reduction requires the Routhian $L-\mu_a\dot x^a$, and the naive substitution of $\dot x(\mu)$ into $L$ gives wrong equations (Remark~\ref{rem:pitfall1}, Example~\ref{ex:kepler}).
\item \emph{Hamiltonian picture.} A conserved momentum is a condition on \emph{points of $T^*Q$}; for abelian $G$ reduction is trivial: set $p_a=\mu_a$ in $H$. A configuration constraint $f=0$ cuts only half of the fibre variables; it must be completed by the momentum condition $\varphi=0$, the pair is second class, and the correct reduction requires the Dirac bracket. Restricting $H$ to $\{f=0\}$ is wrong (Remark~\ref{rem:naiverest}).
\end{itemize}
Conditions on positions are natural on the Lagrangian side, conditions on momenta are natural on the Hamiltonian side, and the Legendre transform exchanges the two. The two reductions therefore look dual in the two formalisms, even though their geometric content is not dual (Section~\ref{sec:unified}).

\subsection{Failure modes}

Both reductions need a regularity (constant-rank) assumption.
\begin{itemize}
\item \emph{Constraints:} the $\dd f_a$ must be independent (otherwise $\mathsf G$ in \eqref{eq:holoPB} is singular and the constraints are not second class); more generally $\rank\Delta$ must be constant on $C$. When it jumps, $\iota_C^*\omega$ has non-constant rank and Theorem~\ref{thm:presymp} does not apply.
\item \emph{Symmetries:} the $G_\mu$-action must be free (or locally free) and $\mu$ regular. Otherwise the quotient is a stratified space (Section~\ref{sec:singular}).
\end{itemize}

\subsection{Nonholonomic constraints}

For velocity constraints that are not integrable, the constrained dynamics is not Hamiltonian in the symplectic sense (Section~\ref{sec:nonhol}), so the left-hand column of Table~\ref{tab:summary} has no symplectic counterpart; and symmetry no longer yields conservation of the full momentum map. We return to this in Section~\ref{sec:nonholsym}.

\begin{small}
\renewcommand{\arraystretch}{1.25}
\setlength{\tabcolsep}{4pt}
\begin{longtable}{@{}>{\raggedright\arraybackslash}p{2.5cm}>{\raggedright\arraybackslash}p{5.5cm}>{\raggedright\arraybackslash}p{5.5cm}@{}}
\caption{Summary: reduction by constraints versus reduction by symmetry.}\label{tab:summary}\\
\toprule
\textbf{Aspect} & \textbf{Constraints (holonomic)} & \textbf{Symmetry (momentum map)}\\
\midrule
\endfirsthead
\toprule
\textbf{Aspect} & \textbf{Constraints (holonomic)} & \textbf{Symmetry (momentum map)}\\
\midrule
\endhead
\bottomrule
\endfoot
Data & Submanifold $Q_c=f^{-1}(0)$ (and, in $T^*Q$, the kinetic metric) & Group action $G$ leaving $L$ or $H$ invariant\\
Logical status & Modelling assumption; reaction forces external to $L$; depends on physical realization (stiff limit) & Consequence of the dynamics (Noether); $\mu$ labels a sector\\
Effect on solutions & $X=X_H+F_C$ with $F_C\ne0$: not solutions of the unreduced system & $F_C=0$: reduced solutions are projections of unreduced solutions\\
Constraint set in $M$ & $S=\{f=0,\ \varphi=0\}$ & $J^{-1}(\mu)$\\
Dirac class & Second class & $\dim G_\mu$ first class $+$ $\dim\mathcal O_\mu$ second class (abelian: all first class)\\
$\ker\iota_C^*\omega$ & $0$ & $T(G_\mu\cdot z)$\\
Operation & Restriction & Restriction, then quotient by $G_\mu$\\
$\dim$ reduced space & $2n-2k$ & $2n-\dim G-\dim G_\mu$ (abelian: $2n-2\dim G$)\\
Reduced symplectic form & Canonical form of $T^*Q_c$ & Canonical $+$ magnetic term depending on $\mu$\\
Parameters & None & $\mu\in\frakg^*$; family $M_\mu$\\
Lost data & Normal variables (fixed by hypothesis); force recovered algebraically & Group variables recovered by quadrature, with geometric phase\\
Lagrangian side & Restrict $L$ to $TQ_c$ (trivial) & Routhian $L-\mu\dot x$ (naive substitution is wrong)\\
Hamiltonian side & Dirac bracket / $T^*Q_c$ (naive restriction is wrong) & Set $p=\mu$ (abelian) and quotient\\
Breakdown & Dependent constraints; rank $\Delta$ jumps & Non-free action, non-regular $\mu$ (stratified spaces)\\
Nonholonomic version & d'Alembert; almost-Poisson; no symplectic reduction & Noether modified; nonholonomic momentum equation; Chaplygin reduction\\
\end{longtable}
\end{small}

%======================================================================
\section{Where the two notions meet}\label{sec:interplay}
%======================================================================

\subsection{Gauge systems, parametrization invariance and energy surfaces}

First-class constraints are the place where the two notions merge mathematically. The standard example is a \emph{parametrized} system: extend $M$ to $\tilde M=M\times T^*\R$ with coordinates $(z,t,p_t)$, $\tilde\omega=\omega+\dd t\wedge\dd p_t$, and impose the single constraint $\chi=p_t+H=0$. It is first class (a single constraint) and $X_\chi=X_H+\partial_t$ so the characteristic leaves of $C=\{\chi=0\}$ are the space-time trajectories; $C/K$ is the space of motions, symplectomorphic to $(M,\omega)$ by evaluation at $t=0$. Reduction by the first-class constraint $\chi$ is a reduction by \emph{time reparametrizations}, a gauge freedom.

Alternatively, let $G=\R$ act on $M$ by the flow of $X_H$ with momentum map $H$ itself. Marsden--Weinstein reduction at $H=E$ gives $H^{-1}(E)/\R$, the space of unparametrized orbits of energy $E$ of dimension $\dim M-2$ (assuming no equilibria and a free proper action). For mechanical $L$ these orbits are the geodesics of the Jacobi metric $(E-V)g$, i.e.\ Jacobi's form of the principle of least action \cite{arnold1989,lanczos1970}. Here a symmetry-based reduction and a gauge-based one coincide.

Even so, the physical status is different. In gauge theories (and parametrized systems) the constraint value is fixed by the theory (Gauss's law $\nabla\cdot E=0$ is the vanishing of the momentum map of the gauge group \cite{henneaux1992,marsden1992}) and states on the same gauge orbit are physically identical: the unreduced phase space is redundant. For a physical symmetry such as rotations, all values of $\mu$ are allowed and points on the same $G$-orbit are physically distinct states (a rotated configuration is a different configuration); the symmetry relates distinct states with the same dynamics instead of identifying them. Noether's first theorem concerns the second situation, Noether's second theorem the first \cite{noether1918}. The mathematics of Theorem~\ref{thm:presymp} does not distinguish them; the physics does.

\subsection{Nonholonomic systems with symmetry}\label{sec:nonholsym}

When constraints are nonholonomic and the system also has symmetries the two reductions interact nontrivially \cite{bkmm1996,bloch2003,cortes2002,marle1995,bates1993}. Let $G$ act on $Q$, with $L$ and the distribution $\mathcal D$ invariant. At $q\in Q$ let $\frakg^q=\{\xi\in\frakg:\xi_Q(q)\in\mathcal D_q\}$ be the symmetry directions allowed by the constraint, and define the \emph{nonholonomic momentum} $J^{\rm nh}(q,\dot q)(\xi)=\pair{\partial L/\partial\dot q}{\xi_Q(q)}$ for $\xi\in\frakg^q$. Then, along solutions of \eqref{eq:nonholeq} and for any curve $\xi(t)\in\frakg^{q(t)}$,
\begin{equation}\label{eq:nhmom}
\frac{\dd}{\dd t}\,J^{\rm nh}(q,\dot q)\big(\xi(t)\big)=\Big\langle\frac{\partial L}{\partial\dot q},\ \big(\dot\xi(t)\big)_Q(q)\Big\rangle,
\end{equation}
the \emph{nonholonomic momentum equation} of Bloch--Krishnaprasad--Marsden--Murray \cite{bkmm1996}. Thus the Noether conservation law survives only for symmetries $\xi$ with $\xi_Q\in\mathcal D$ everywhere (and $\xi$ constant) \cite{bates1993,sniatycki1998}; the momenta of symmetries transverse to the constraint distribution are generally not conserved, because the constraint force has a component along $\xi_Q$. In the favourable case where $\mathcal D$ is a horizontal space of a principal connection on $Q\to Q/G$ (Chaplygin systems) the symmetry-reduced equations live on $T(Q/G)$, are of Lagrange--d'Alembert type, and contain a gyroscopic term built from the curvature of the constraint connection; they are in general not Lagrangian, and sometimes become Hamiltonian after a time reparametrization (Chaplygin's reducing multiplier) \cite{chaplygin1911,koiller1992,ehlers2005}. Thus in nonholonomic mechanics a symmetry reduction can \emph{destroy} Hamiltonian structure rather than preserve it, in contrast with Theorem~\ref{thm:MWM}.

%======================================================================
\section{Conclusion}\label{sec:conclusion}
%======================================================================

We return to the questions of the introduction.

\begin{enumerate}[label=(Q\arabic*)]
\item \emph{What is reduced, and how?} In both cases the symplectic manifold $(M,\omega)$ is replaced by $(\bar C,\bar\omega)$, obtained from a submanifold $C$ by dividing out the kernel of $\iota_C^*\omega$. For a holonomic constraint the kernel vanishes and the operation is a restriction (Dirac's second-class reduction, $T^*Q_c$). For a symmetry the submanifold is a momentum level set and the kernel is the orbit of the isotropy group, so the operation is a restriction and a quotient (Marsden--Weinstein--Meyer; Routh in the abelian case).
\item \emph{Consequence or modification?} Symmetry reduction is a consequence of the dynamics: the constraint force is identically zero and reduced trajectories are projections of genuine trajectories. Constraint reduction is a modification: a reaction force acts, and the reduced trajectories are not trajectories of the unconstrained system.
\item \emph{Is one a special case of the other?} Mathematically, symmetry reduction is a case of Dirac's theory of constrained Hamiltonian systems, with constraints $J-\mu$ of which $\dim G_\mu$ are first class and $\dim\mathcal O_\mu$ second class (Proposition~\ref{prop:symclass}); the dual statement, that a holonomic constraint is a symmetry reduction, is false. Physically they differ: a constraint is imposed, a symmetry is exploited; and in the Lagrangian and Hamiltonian formalisms the roles of ``trivial'' and ``subtle'' are interchanged. The two meet in gauge systems, parametrized systems and energy-surface reduction, and they interact in subtle ways for nonholonomic constraints.
\end{enumerate}

Practical rule of thumb: \emph{if you are about to ``eliminate a variable using a conservation law'', you are doing symmetry reduction and no new force is introduced; if you are ``eliminating a variable using a relation that you were given'', you are doing constraint reduction and the equations of motion have been changed.}

For further reading on the textbook level we recommend \cite{goldstein2002,landau1976,lanczos1970,arnold1989,josesaletan1998}; for symplectic geometry \cite{guilleminsternberg1984,libermannmarle1987,souriau1970,cannas2001}; for constrained Hamiltonian dynamics \cite{sundermeyer1982,henneaux1992,gnh1978}; for the full theory of momentum maps and reduction \cite{am1978,marsden1992,mr1999,op2004,mmopr2007}.

%======================================================================
\begin{thebibliography}{99}
%======================================================================
\small
\bibitem{am1978} R. Abraham and J. E. Marsden, \emph{Foundations of Mechanics}, 2nd ed., Benjamin/Cummings, Reading, MA, 1978.
\bibitem{andersonbergmann1951} J. L. Anderson and P. G. Bergmann, Constraints in covariant field theories, \emph{Phys. Rev.} \textbf{83} (1951), 1018--1025.
\bibitem{arnold1989} V. I. Arnold, \emph{Mathematical Methods of Classical Mechanics}, 2nd ed., Graduate Texts in Math.\ 60, Springer, New York, 1989.
\bibitem{akn} V. I. Arnold, V. V. Kozlov and A. I. Neishtadt, \emph{Mathematical Aspects of Classical and Celestial Mechanics}, 3rd ed., Encyclopaedia Math.\ Sci.\ 3, Springer, Berlin, 2006.
\bibitem{bates1993} L. Bates and J. \'Sniatycki, Nonholonomic reduction, \emph{Rep. Math. Phys.} \textbf{32} (1993), 99--115.
\bibitem{bergmann1949} P. G. Bergmann, Non-linear field theories, \emph{Phys. Rev.} \textbf{75} (1949), 680--685.
\bibitem{bkmm1996} A. M. Bloch, P. S. Krishnaprasad, J. E. Marsden and R. M. Murray, Nonholonomic mechanical systems with symmetry, \emph{Arch. Rational Mech. Anal.} \textbf{136} (1996), 21--99.
\bibitem{bloch2003} A. M. Bloch, \emph{Nonholonomic Mechanics and Control}, Interdisciplinary Applied Math.\ 24, Springer, New York, 2003.
\bibitem{bornemann1997} F. A. Bornemann and C. Sch\"utte, Homogenization of Hamiltonian systems with a strong constraining potential, \emph{Physica D} \textbf{102} (1997), 57--77.
\bibitem{cannas2001} A. Cannas da Silva, \emph{Lectures on Symplectic Geometry}, Lecture Notes in Math.\ 1764, Springer, Berlin, 2001.
\bibitem{cmr2001} H. Cendra, J. E. Marsden and T. S. Ratiu, Lagrangian reduction by stages, \emph{Mem. Amer. Math. Soc.} \textbf{152} (2001), no.~722.
\bibitem{chaplygin1911} S. A. Chaplygin, On the theory of motion of nonholonomic systems. The reducing-multiplier theorem, \emph{Regul. Chaotic Dyn.} \textbf{13} (2008), 369--376 (translation of the 1911 Russian original).
\bibitem{cortes2002} J. Cort\'es Monforte, \emph{Geometric, Control and Numerical Aspects of Nonholonomic Systems}, Lecture Notes in Math.\ 1793, Springer, Berlin, 2002.
\bibitem{dacosta1981} R. C. T. da Costa, Quantum mechanics of a constrained particle, \emph{Phys. Rev. A} \textbf{23} (1981), 1982--1987.
\bibitem{dalembert} J. le Rond d'Alembert, \emph{Trait\'e de dynamique}, David, Paris, 1743.
\bibitem{dirac1950} P. A. M. Dirac, Generalized Hamiltonian dynamics, \emph{Canad. J. Math.} \textbf{2} (1950), 129--148.
\bibitem{dirac1964} P. A. M. Dirac, \emph{Lectures on Quantum Mechanics}, Belfer Graduate School of Science, Yeshiva University, New York, 1964 (Dover reprint, 2001).
\bibitem{ehlers2005} K. Ehlers, J. Koiller, R. Montgomery and P. M. Rios, Nonholonomic systems via moving frames: Cartan equivalence and Chaplygin Hamiltonization, in \emph{The Breadth of Symplectic and Poisson Geometry}, Progr. Math.\ 232, Birkh\"auser, Boston, 2005, pp.~75--120.
\bibitem{goldstein2002} H. Goldstein, C. Poole and J. Safko, \emph{Classical Mechanics}, 3rd ed., Addison-Wesley, San Francisco, 2002.
\bibitem{gotay1982} M. J. Gotay, On coisotropic imbeddings of presymplectic manifolds, \emph{Proc. Amer. Math. Soc.} \textbf{84} (1982), 111--114.
\bibitem{gotaynester1979} M. J. Gotay and J. M. Nester, Presymplectic Lagrangian systems I: the constraint algorithm and the equivalence theorem, \emph{Ann. Inst. H. Poincar\'e Sect. A} \textbf{30} (1979), 129--142.
\bibitem{gnh1978} M. J. Gotay, J. M. Nester and G. Hinds, Presymplectic manifolds and the Dirac--Bergmann theory of constraints, \emph{J. Math. Phys.} \textbf{19} (1978), 2388--2399.
\bibitem{guilleminsternberg1984} V. Guillemin and S. Sternberg, \emph{Symplectic Techniques in Physics}, Cambridge University Press, Cambridge, 1984.
\bibitem{henneaux1992} M. Henneaux and C. Teitelboim, \emph{Quantization of Gauge Systems}, Princeton University Press, Princeton, 1992.
\bibitem{hmr1998} D. D. Holm, J. E. Marsden and T. S. Ratiu, The Euler--Poincar\'e equations and semidirect products with applications to continuum theories, \emph{Adv. Math.} \textbf{137} (1998), 1--81.
\bibitem{josesaletan1998} J. V. Jos\'e and E. J. Saletan, \emph{Classical Dynamics: A Contemporary Approach}, Cambridge University Press, Cambridge, 1998.
\bibitem{koiller1992} J. Koiller, Reduction of some classical non-holonomic systems with symmetry, \emph{Arch. Rational Mech. Anal.} \textbf{118} (1992), 113--148.
\bibitem{kummer1981} M. Kummer, On the construction of the reduced phase space of a Hamiltonian system with symmetry, \emph{Indiana Univ. Math. J.} \textbf{30} (1981), 281--291.
\bibitem{lagrange} J.-L. Lagrange, \emph{M\'ecanique analytique}, Paris, 1788.
\bibitem{lanczos1970} C. Lanczos, \emph{The Variational Principles of Mechanics}, 4th ed., University of Toronto Press, Toronto, 1970 (Dover reprint, 1986).
\bibitem{landau1976} L. D. Landau and E. M. Lifshitz, \emph{Mechanics}, 3rd ed., Pergamon, Oxford, 1976.
\bibitem{libermannmarle1987} P. Libermann and C.-M. Marle, \emph{Symplectic Geometry and Analytical Mechanics}, Reidel, Dordrecht, 1987.
\bibitem{marle1995} C.-M. Marle, Reduction of constrained mechanical systems and stability of relative equilibria, \emph{Comm. Math. Phys.} \textbf{174} (1995), 295--318.
\bibitem{marsden1992} J. E. Marsden, \emph{Lectures on Mechanics}, London Math.\ Soc.\ Lecture Note Ser.\ 174, Cambridge University Press, Cambridge, 1992.
\bibitem{mmr1990} J. E. Marsden, R. Montgomery and T. S. Ratiu, Reduction, symmetry, and phases in mechanics, \emph{Mem. Amer. Math. Soc.} \textbf{88} (1990), no.~436.
\bibitem{mmopr2007} J. E. Marsden, G. Misio{\l}ek, J.-P. Ortega, M. Perlmutter and T. S. Ratiu, \emph{Hamiltonian Reduction by Stages}, Lecture Notes in Math.\ 1913, Springer, Berlin, 2007.
\bibitem{marsdenperlmutter2000} J. E. Marsden and M. Perlmutter, The orbit bundle picture of cotangent bundle reduction, \emph{C. R. Math. Acad. Sci. Soc. R. Can.} \textbf{22} (2000), 33--54.
\bibitem{mr1999} J. E. Marsden and T. S. Ratiu, \emph{Introduction to Mechanics and Symmetry}, 2nd ed., Texts in Applied Math.\ 17, Springer, New York, 1999.
\bibitem{mr1986} J. E. Marsden and T. Ratiu, Reduction of Poisson manifolds, \emph{Lett. Math. Phys.} \textbf{11} (1986), 161--169.
\bibitem{mrs2000} J. E. Marsden, T. S. Ratiu and J. Scheurle, Reduction theory and the Lagrange--Routh equations, \emph{J. Math. Phys.} \textbf{41} (2000), 3379--3429.
\bibitem{marsdenscheurle1993} J. E. Marsden and J. Scheurle, Lagrangian reduction and the double spherical pendulum, \emph{Z. Angew. Math. Phys.} \textbf{44} (1993), 17--43.
\bibitem{mw1974} J. E. Marsden and A. Weinstein, Reduction of symplectic manifolds with symmetry, \emph{Rep. Math. Phys.} \textbf{5} (1974), 121--130.
\bibitem{meyer1973} K. R. Meyer, Symmetries and integrals in mechanics, in \emph{Dynamical Systems} (Proc. Sympos., Univ. Bahia, Salvador, 1971), M. Peixoto, ed., Academic Press, New York, 1973, pp.~259--272.
\bibitem{montgomery1991} R. Montgomery, How much does the rigid body rotate? A Berry's phase from the 18th century, \emph{Amer. J. Phys.} \textbf{59} (1991), 394--398.
\bibitem{noether1918} E. Noether, Invariante Variationsprobleme, \emph{Nachr. Ges. Wiss. G\"ottingen, Math.-Phys. Kl.} (1918), 235--257.
\bibitem{op2004} J.-P. Ortega and T. S. Ratiu, \emph{Momentum Maps and Hamiltonian Reduction}, Progr. Math.\ 222, Birkh\"auser, Boston, 2004.
\bibitem{poincare1901} H. Poincar\'e, Sur une forme nouvelle des \'equations de la m\'ecanique, \emph{C. R. Acad. Sci. Paris} \textbf{132} (1901), 369--371.
\bibitem{routh1877} E. J. Routh, \emph{A Treatise on the Stability of a Given State of Motion}, Macmillan, London, 1877.
\bibitem{rubin1957} H. Rubin and P. Ungar, Motion under a strong constraining force, \emph{Comm. Pure Appl. Math.} \textbf{10} (1957), 65--87.
\bibitem{satzer1977} W. J. Satzer, Jr., Canonical reduction of mechanical systems invariant under abelian group actions with an application to celestial mechanics, \emph{Indiana Univ. Math. J.} \textbf{26} (1977), 951--976.
\bibitem{sjamaarlerman1991} R. Sjamaar and E. Lerman, Stratified symplectic spaces and reduction, \emph{Ann. of Math.} \textbf{134} (1991), 375--422.
\bibitem{smale1970} S. Smale, Topology and mechanics I, II, \emph{Invent. Math.} \textbf{10} (1970), 305--331; \textbf{11} (1970), 45--64.
\bibitem{sniatycki1998} J. \'Sniatycki, Nonholonomic Noether theorem and reduction of symmetries, \emph{Rep. Math. Phys.} \textbf{42} (1998), 5--23.
\bibitem{souriau1970} J.-M. Souriau, \emph{Structure des syst\`emes dynamiques}, Dunod, Paris, 1970.
\bibitem{sundermeyer1982} K. Sundermeyer, \emph{Constrained Dynamics}, Lecture Notes in Phys.\ 169, Springer, Berlin, 1982.
\bibitem{takens1980} F. Takens, Motion under the influence of a strong constraining force, in \emph{Global Theory of Dynamical Systems}, Lecture Notes in Math.\ 819, Springer, Berlin, 1980, pp.~425--445.
\bibitem{vankampen1984} N. G. van Kampen and J. J. Lodder, Constraints, \emph{Amer. J. Phys.} \textbf{52} (1984), 419--424.
\bibitem{vanderschaft1994} A. J. van der Schaft and B. M. Maschke, On the Hamiltonian formulation of nonholonomic mechanical systems, \emph{Rep. Math. Phys.} \textbf{34} (1994), 225--233.
\end{thebibliography}

\end{document}
